\chapter{Numerische Rekonstruktion für einen prüfbaren Kandidaten}

\section{Warum eine Krümmungsänderung Numerik benötigt}

Im durchgehenden Beispiel verändern die freien Variablen ausgewählte
Koeffizienten eines intrinsischen Krümmungsgesetzes. Der feste Endpunkt, die
Endrichtung und die Vermessungsbeobachtungen sind jedoch räumliche Größen.
Jeder Versuchsvektor \(x\) muss deshalb rekonstruiert werden, bevor das Problem
seine Übereinstimmung bewerten kann.

Sei
\[
\kappa_x:[0,L]\rightarrow\R
\]
die Kandidatenkrümmung und \((\gamma_0,\theta_0)\) die feste Anfangspose. Das
Modell \(M\) erzeugt
\[
\theta_x'(s)=\kappa_x(s),\qquad
\gamma_x'(s)=(\cos\theta_x(s),\sin\theta_x(s)).
\]
Diese Rekonstruktion erzeugt eine für die Evaluation verwendete Darstellung.
Sie ändert weder die Alignment-Identität noch macht sie die dargestellte Kurve
zu einem angenommenen Engineering-Zustand.

\section{Vom kontinuierlichen Modell zu berechenbaren Abtastwerten}

\paragraph{Lesebrücke.}
Die Gleichungen beschreiben die Abhängigkeit exakt; ein Solver muss sie jedoch
endlich und wiederholt auswerten. Diskretisierung ist die erklärte
Approximation, die dies ermöglicht.

Wähle
\[
0=s_0<s_1<\dots<s_N=L,\qquad \Delta s_i=s_{i+1}-s_i .
\]
Ein Mittelpunktsschritt liefert beispielsweise
\[
\theta_{i+\frac12}=\theta_i+
 \frac12\kappa_x(s_i)\Delta s_i,
\]
\[
\gamma_{i+1}=\gamma_i+\Delta s_i
(\cos\theta_{i+\frac12},\sin\theta_{i+\frac12}),\qquad
\theta_{i+1}=\theta_i+\kappa_x(s_i)\Delta s_i.
\]
Andere Integrationsverfahren können geeignet sein. Verfahren, Toleranz und
Abtastregel gehören zu \(M\); sie sind kein unsichtbares
Implementierungsdetail.

\section{Residuen beantworten erklärte Fragen}

Für die feste Zielendpose \((\bar\gamma_L,\bar\theta_L)\) kann ein Modell das
Endresiduum
\[
r_{\mathrm{end}}(x)=
\begin{bmatrix}
\gamma_x(L)-\bar\gamma_L\\
\operatorname{wrap}(\theta_x(L)-\bar\theta_L)
\end{bmatrix}
\]
erzeugen. Für Vermessungsbeobachtungen \(y_j\) kann ein Beobachtungsmodell
\(h_j\) liefern:
\[
r_{\mathrm{obs},j}(x)=h_j(\gamma_x)-y_j.
\]
Diese Residuen beantworten verschiedene Fragen und können verschiedene
Einheiten, Unsicherheiten und Verantwortlichkeiten besitzen. Ein
Nebenbedingungsresiduum beschreibt Zulässigkeit; ein Beobachtungsresiduum
beschreibt eine Abweichung unter einem Beobachtungsmodell. Sie dürfen nicht
stillschweigend zusammengeführt werden.

Gehen Residuenfamilien in ein gemeinsames Kleinste-Quadrate-Ziel ein, muss die
Skalierung erklärt werden, etwa
\[
J(x)=
\frac12\|W_{\mathrm{end}}r_{\mathrm{end}}(x)\|_2^2+
\frac12\sum_j
\|W_{\mathrm{obs},j}r_{\mathrm{obs},j}(x)\|_2^2+
J_{\mathrm{reg}}(x).
\]
Die Matrizen \(W\) machen Einheiten und relativen Einfluss sichtbar. Ihre
Werte stammen aus erklärtem Zweck, Unsicherheitsmodell oder verantwortlicher
ingenieurtechnischer Festlegung, nicht aus dem numerischen Verfahren.

\section{Genauigkeit gehört zur Ergebnisevidenz}

\paragraph{Lesebrücke.}
Ein Kandidat kann wegen eines groben Gitters zulässig erscheinen, obwohl das
zugrunde liegende kontinuierliche Modell die Nebenbedingung verletzt.
Numerischer Fehler muss deshalb mit dem Ergebnis weitergegeben werden.

Relevante Beiträge sind:
\begin{itemize}
\item Integrations- und Diskretisierungsfehler;
\item Interpolations- oder Quadraturfehler in Beobachtungsresiduen;
\item Gleitkomma- und lineare Lösungsfehler;
\item Empfindlichkeit der rekonstruierten Lage gegenüber
      Krümmungsstörungen;
\item Modellform- und Beobachtungsunsicherheit, die numerische Verfeinerung
      nicht beseitigen kann.
\end{itemize}

Ein belastbarer Lauf dokumentiert Verfahren und Gitter, Toleranzen,
Verfeinerungsprüfungen sowie auf einem unabhängigen oder feineren
Bewertungsgitter neu berechnete Residuen. Ein gleitender Ursprung oder lokaler
metrischer Rahmen kann numerische Verluste bei großen Koordinatenwerten
verringern; diese Darstellungsänderung verändert das intrinsische Alignment
nicht.

\section{Warum als Nächstes Ableitungen und dünne Besetzung erscheinen}

Wiederholte Rekonstruktion macht die Abhängigkeit
\[
x\longmapsto\gamma_x\longmapsto r(x)\longmapsto J(x)
\]
rechnerisch sichtbar. Ableitungen beschreiben, wie eine erlaubte
Parameteränderung die Residuen beeinflusst. Dünne Besetzung hält fest, dass
ein lokaler Krümmungsparameter gewöhnlich nur strukturierte Teile des
zusammengesetzten Problems beeinflusst. Beides beantwortet Anforderungen an
Aufwand und Nachvollziehbarkeit; keines wählt das ingenieurtechnische Ziel.

Das Numerikkapitel hat damit den für das Optimierungsproblem erforderlichen
Evaluator geliefert:
\[
(x;M,q)\longmapsto
\bigl(\widehat\gamma_x,r(x),J(x),\varepsilon_{\mathrm{num}}\bigr).
\]
Sein Ergebnis ist Evidenz über einen Versuch, keine Entscheidung über das
Alignment.

\begin{summary}
Numerische Rekonstruktion macht eine erklärte Krümmungsänderung auswertbar.
Diskretisierung, Skalierung, Toleranzen und Fehlerschätzungen bleiben sichtbar,
damit ein scheinbar gutes Solver-Ergebnis geprüft werden kann. Das nächste
Kapitel kann nun untersuchen, wie die Freiheitsgrade effizient durchsucht
werden, ohne die Verantwortung für Problem oder Entscheidung zu verschieben.
\end{summary}
